% !TeX program = lualatex
% Compile twice: lualatex 82.tex
% Standalone research notebook; original section preserved.
% Research additions: 27 September 2026. No external TeX input or BibTeX file.
\documentclass[11pt,a4paper]{article}
\usepackage[margin=24mm,headheight=14pt]{geometry}
\usepackage{fontspec}
\setmainfont{Latin Modern Roman}
\setsansfont{Latin Modern Sans}
\setmonofont{Latin Modern Mono}
\usepackage{amsmath,amssymb,mathtools}
\usepackage{unicode-math}
\setmathfont{Latin Modern Math}
\usepackage{microtype}
\usepackage{enumitem,xurl,needspace}
\usepackage[dvipsnames]{xcolor}
\usepackage{hyperref}
\hypersetup{unicode=true,colorlinks=true,linkcolor=MidnightBlue,urlcolor=MidnightBlue,
 pdftitle={Research notebook - Problem 82},pdfauthor={Open Problems Math},bookmarksnumbered=true}
\usepackage{bookmark}
\usepackage{titlesec}
\titleformat{\section}{\Large\bfseries\raggedright}{\thesection}{0.7em}{}
\usepackage{fancyhdr}
\pagestyle{fancy}\fancyhf{}
\fancyhead[L]{\small\sffamily Open Problems Math | Research notebook}
\fancyhead[R]{\small\sffamily Problem 82 | 27 September 2026}
\fancyfoot[C]{\thepage}
\setlength{\parindent}{0pt}
\setlength{\parskip}{5pt plus 1pt}
\setlength{\emergencystretch}{3em}
\setcounter{tocdepth}{1}
\setcounter{secnumdepth}{1}
\setlist[itemize]{leftmargin=3em,itemsep=5pt,topsep=4pt,parsep=0pt}
\allowdisplaybreaks
\numberwithin{equation}{section}
\newcommand{\N}{\mathbb{N}}
\newcommand{\Z}{\mathbb{Z}}
\newcommand{\Q}{\mathbb{Q}}
\newcommand{\R}{\mathbb{R}}
\newcommand{\C}{\mathbb{C}}
\newcommand{\F}{\mathbb{F}}
\newcommand{\A}{\mathbb{A}}
\newcommand{\PP}{\mathbb P}
\newcommand{\E}{\mathbb{E}}
\newcommand{\eps}{\varepsilon}
\newcommand{\dd}{\,\mathrm d}
\newcommand{\sref}[2]{\hyperlink{source:#1:#2}{[S#2]}}
\newcommand{\eref}[2]{\hyperlink{extra:#1:#2}{[E#2]}}
\newif\ifshowcatalogue
\showcataloguetrue
\newcommand{\cataloguescope}[1]{\ifshowcatalogue
 \paragraph{Exact scope in the supplied catalogue.}
 \begin{quote}\small #1\end{quote}\fi}
\begin{document}
\setcounter{section}{81}
\section{\texorpdfstring{Zilber-Pink Conjecture on Unlikely Intersections}{Zilber-Pink Conjecture on Unlikely Intersections}}\label{82}
\subsection{Definitions and mathematical statement}
A special subvariety of a mixed Shimura variety is one arising from the appropriate Shimura subdatum and Hecke-translate construction. For irreducible $W\subseteq S$, the expected intersection dimension with $T$ is $\dim W+\dim T-\dim S$. The record calls an irreducible $A\subseteq W\cap T$ atypical if
\[
 \dim A>\dim W+\dim T-\dim S.
\]
The claim is that only finitely many such $A$ are maximal under inclusion. Equivalently their union is the union of finitely many atypical subvarieties. The source's special-subvariety convention is essential; replacing ``special'' with an arbitrary algebraic family would give a different problem.
\par\smallskip\textit{Formulation and concept references:} \sref{82}{1}.
\cataloguescope{Let S be a mixed Shimura variety and W an irreducible subvariety of S defined over C. Call a subvariety A of W atypical if there is a special subvariety T of S with A contained in the intersection of T and W and dim A > dim W + dim T - dim S. Prove that the union of all atypical subvarieties of W is a finite union of atypical subvarieties, equivalently that W contains only finitely many maximal atypical subvarieties.}
\subsection{Short English statement}
Are all unexpectedly large special intersections inside a fixed Shimura subvariety controlled by finitely many maximal exceptional pieces?
% BEGIN RESEARCH_ADDITION ATTEMPT_82
\subsection{Research attempt}

\paragraph{Current frontier and special-subvariety convention.}
Pila's Section~2.3 was inspected for the mixed Shimura formulation and its
countable collection of special subvarieties \eref{82}{1}. The full
complex-field statement must not be replaced by a number-field case.
The work on varieties not defined over $\overline\Q$ in \sref{82}{2}
addresses that distinction, but its particular reduction and hypotheses
cannot simply be declared a theorem for every mixed Shimura variety.
The counting ingredient used below is the Pila--Wilkie bound on the
transcendental part of a fixed definable set \eref{82}{2}.

\paragraph{Attack route.}
A tempting qualitative route says that the atypical union is algebraic by
Noetherianity. This is false without a closedness theorem. An arithmetic
route compares large collections of conjugates with subpolynomial definable
point counts, but needs bounded-complexity coding and control of
semialgebraic blocks. The attempt proves a precise closedness lemma and
an explicit counting criterion, then identifies the missing geometric
bridges for the exact all-complex, mixed setting.

\paragraph{Attempt: closedness would be enough, but is not automatic.}
Write $U\subseteq W$ for the union of all atypical subvarieties.
There are only countably many special $T$. Each intersection $W\cap T$
has finitely many irreducible components. If an irreducible
$A\subseteq W\cap T$ is atypical, an irreducible component $B$ of that
intersection containing it satisfies $\dim B\geq\dim A$ and the same
strict dimension inequality, hence is itself atypical. Therefore $U$ is
a countable union $\bigcup_i B_i$ of closed irreducible atypical subvarieties.
The countability here is for special, not arbitrary weakly special,
subvarieties \eref{82}{1}.

\emph{Lemma.} For such a countable family over $\C$, if $U$ is Zariski
closed, it is a finite union of atypical subvarieties.
To prove it, write the closed set $U$ as a finite union of its irreducible
components $Z_1,\ldots,Z_s$. An irreducible complex variety cannot be a
countable union of proper closed subsets. One proof takes an affine open,
a finite Noether normalization onto affine space, and images of the proposed
proper closed subsets. These images are proper closed sets of smaller
dimension. All their defining coefficients lie in one countable subfield
of $\C$. A point of affine space whose coordinates are algebraically
independent over that field avoids all of them; such coordinates can be
chosen successively because each algebraic closure of a countable field
is countable. A point in the finite nonempty fibre of the normalization
then avoids the proposed cover. Dimension zero is immediate.

Apply this fact to
$Z_j=\bigcup_i(Z_j\cap B_i)$. Some $B_i$ contains $Z_j$. Since
$B_i\subseteq U$ and $Z_j$ is a maximal irreducible component of $U$,
$B_i=Z_j$. Thus every $Z_j$ is atypical, proving the lemma. The converse,
that a finite union is closed, is immediate. This supplies an exact
closedness reformulation in this setting.

The missing statement is precisely that $U$ is closed. Noetherianity
controls descending chains of closed sets and finite irreducible
decompositions of a \emph{closed} set, not arbitrary unions. An infinite
countable set of complex points in $\A^1$ is already a counterexample to
the naive inference: every point is closed, but their union is not closed.
Taking the closure of $U$ can add subvarieties not yet shown atypical,
which is the issue the conjecture requires one to solve.

\paragraph{Attempt: a finite-complexity counting criterion.}
Consider a collection of candidate maximal atypical objects, possibly
subvarieties encoded by parameters, with complexity $\Delta\geq1$.
Suppose one can prove all the following for this fixed $W$: there is a
single definable set $Z\subseteq\R^m$; every candidate yields at least
$c\Delta^\delta$ distinct rational codes in $Z^{\rm tr}$; their heights
are at most $C\Delta^b$; and the relation from codes to candidates has
finite fibres. Here $b,\delta,c,C>0$ are fixed, and $Z^{\rm tr}$ excludes
the positive-dimensional connected semialgebraic subsets.
The Pila--Wilkie theorem states
\[
 \#\{z\in Z^{\rm tr}\cap\Q^m:H(z)\leq H\}\ll_{Z,\varepsilon}H^\varepsilon
 \qquad(\varepsilon>0)
\]
\eref{82}{2}. Therefore each candidate would satisfy
\[
 c\Delta^\delta\leq C_{Z,\varepsilon}(C\Delta^b)^\varepsilon.
\]
Choose $\varepsilon=\delta/(2b)$. It follows that $\Delta$ is bounded
uniformly over the candidates. The height bound then places every candidate
above at least one member of a finite set of rational codes. Finite fibres
make the candidate collection finite.

This is a rigorous conditional counting argument, including a finiteness
condition sometimes lost after obtaining bounded complexity. A bound on
complexity alone is not useful if infinitely many candidates may share
that complexity and the same code. Rational codes are an explicitly
sufficient convention here, not an assertion that the actual mixed Shimura
objects already possess such codes. Versions with bounded-degree algebraic
codes need their corresponding counting theorem and degree bounds; neither
is silently substituted.

\paragraph{Attempt: a completely checked toric calibration.}
As a separate toric analogue, let $W\subset(\C^\times)^2$ be $x+y=1$
and let special subvarieties mean torsion cosets. The curve is not contained
in a proper torsion coset: an identity $x^a(1-x)^b=\zeta$ on $W$ would
have zero orders at $x=0$ and $x=1$, forcing $a=b=0$.
Thus positive-dimensional intersections with proper torsion cosets are
not possible. The atypical points in this two-dimensional ambient torus
are its torsion points on $W$.
If both $x$ and $y=1-x$ are roots of unity, $|x|=|1-x|=1$, and hence
\[
 1=|1-x|^2=2-2\operatorname{Re}(x).
\]
So $\operatorname{Re}(x)=1/2$ and
$x=(1\pm i\sqrt3)/2$, with $y$ its conjugate. Both possibilities really
are roots of unity. There are exactly two such points. The calculation
illustrates the need for extra arithmetic rigidity to make a countable
special-point union finite; countability alone did not do it.
It is not offered as a solution for an arbitrary mixed Shimura variety.

\paragraph{Stress test.}
A lower bound for Galois orbits cannot be used without showing that the
relevant conjugates stay in the fixed $W$. For arbitrary $W/\C$, that
stability is not automatic. Spreading out, field-of-definition parameters,
and special closures must be handled rather than replacing $W$ by a
variety over $\overline\Q$. The source's scope explicitly retains these
complex varieties.

Also, definable counting permits many points on its semialgebraic part.
If the hypothetical codes land there, the displayed upper bound does not
apply. One would need a functional-transcendence or block-analysis theorem
turning such families into larger atypical subvarieties and then exploit
maximality. Positive-dimensional maximal subvarieties need their own
parameter spaces, degree control and finite-fibre coding; a theorem only
about isolated special points does not settle their finiteness. None of
those geometric bridges is supplied by the elementary counting inequality.

\paragraph{Outcome.}
\textbf{Outcome: Conditional reduction.}

The attempt proves the closedness-to-finiteness lemma under the exact
countability convention and a quantitative coding criterion that would
force finitely many maximal objects. It checks a small toric analogue and
identifies why Noetherianity and raw Galois-orbit language do not close the
general argument. No new mixed Shimura case is proved. A completion needs
the actual closedness theorem or the missing coding, orbit, and block
control for every $W/\C$ in the original scope; the elementary reductions
carry no priority claim.
% END RESEARCH_ADDITION ATTEMPT_82
\subsection{Sources}
\begin{itemize}\raggedright
\item[\textnormal{[S1]}]\hypertarget{source:82:1}{} J. Pila, 'Combining the conjectures of Schanuel and Zilber-Pink', Rend. Lincei Mat. Appl. 36 (2025), 653-670, DOI 10.4171/RLM/1086; accepted 1 July 2025.
\url{https://ems.press/content/serial-article-files/52570}.
{\small\textit{Catalogue formulation source.}}
\item[\textnormal{[S2]}]\hypertarget{source:82:2}{} The Zilber-Pink conjecture for varieties not defined over Qbar.
\url{https://arxiv.org/html/2504.00865v1}.
{\small\textit{Catalogue research/status reference; not a proof endorsement.}}
\item[\textnormal{[S3]}]\hypertarget{source:82:3}{} A note on unlikely intersections in Shimura varieties.
\url{https://arxiv.org/html/2209.07967}.
{\small\textit{Catalogue research/status reference; not a proof endorsement.}}
% BEGIN RESEARCH_ADDITION REFERENCES_82
\item[\textnormal{[E1]}]\hypertarget{extra:82:1}{} Jonathan Pila,
\emph{Combining the conjectures of Schanuel and Zilber--Pink},
Rendiconti Lincei -- Matematica e Applicazioni 36 (2025), 653--670,
Section 2.3 and the discussion of special subvarieties.
\url{https://ems.press/content/serial-article-files/52570}.
{\small\textit{Exact mixed-Shimura formulation inspected 27 September 2026;
the toric example in this notebook is explicitly only an analogue.}}
\item[\textnormal{[E2]}]\hypertarget{extra:82:2}{} Jonathan Pila and
Alex J. Wilkie, \emph{The rational points of a definable set},
Duke Mathematical Journal 133 (2006), 591--616, main counting theorem.
\url{https://eprints.maths.manchester.ac.uk/942/}.
{\small\textit{Primary institutional manuscript record inspected
27 September 2026; the theorem is used only on a hypothetical fixed
definable transcendental part with rational codes and explicit height bounds.}}
% END RESEARCH_ADDITION REFERENCES_82
\end{itemize}

\end{document}
